\documentclass[10pt,a4paper]{amsart}
\usepackage[margin=19mm]{geometry}
\usepackage{amsmath,amssymb,mathtools}
\usepackage[scr=rsfs,cal=euler]{mathalfa}
\usepackage{xfrac}
\usepackage[unicode,hidelinks,bookmarks=false]{hyperref}
\newcommand{\ie}{i.e.\ }
\newcommand{\cf}{cf.\ }
\newcommand{\resp}{respectively\ }
\newcommand{\Subsection}{\subsection}
\providecommand{\nobreakdash}{\nobreak\hbox{-}}
\setlength{\emergencystretch}{2em}
\multlinegap0pt
\usepackage{bm}
\usepackage{bbm}
\usepackage{dsfont}
\usepackage{xspace}
\usepackage{cancel}
\usepackage{mathtools}
\usepackage{amsthm}
\makeatletter
\@ifundefined{theorem}{\newtheorem{theorem}{Theorem}}{}
\makeatother
\makeatletter
\expandafter\ifx\csname dedication\endcsname\relax
\newenvironment{dedication}
{\vspace{0ex}\begin{quotation}\begin{center}\begin{em}}
\fi
\def\min{{\operatorname{min}\nolimits}}
\makeatother
\title[Big categorification on towers of classical groups and wreat]{Big categorification on towers of classical groups and wreath product groups}
\author{Xin Huang, Pengcheng Li,, Yaolong Shen}
\date{Source: arXiv:2608.17709v3 (CC BY 4.0)}
\begin{document}
\maketitle

\noindent\emph{Source and licence.} Big categorification on towers of classical groups and wreath product groups. Version arXiv:2608.17709v3 (CC BY 4.0), distributed under the Creative Commons Attribution 4.0 International licence, as recorded in the arXiv licence metadata for this exact version and shown on the abstract page https://arxiv.org/abs/2608.17709v3. Full rights notice: licenses/arxiv-2608.17709-CC-BY-4.0.txt.\par\smallskip

\noindent\textit{Editorial scope.} Counted unit: the Theorem whose source label is \texttt{Prop:TypeBC} in the article (source0.tex lines 6794--6822, source label \texttt{Prop:TypeBC}), delivered together with the article's own complete original proof of it (source0.tex lines 6824--6915, one \texttt{proof} environment of 92 lines). No mathematics is written, repaired or re-derived: statement and proof are the article's own text line for line. Required context retained uncounted: none. Every definition, display and label of those lines is retained unchanged. Omitted from this package and not redistributed: 1--6793, 6823--6823, 6916--11089. Nothing inside a retained excerpt is omitted or altered: every retained body is delivered exactly as the article's lines are written, so no omission receipt of any kind accompanies this package. Citations: the retained excerpts carry no citation command at all, so no bibliography entry of the article is transcribed and no reference is redistributed. Reference adaptation: none was needed and none was made; every reference inside the delivered unit resolves inside it, so no unresolved reference or citation is produced. Printed slips kept literal: no printed word, symbol, hypothesis or number of the counted statement or of its proof was altered. Layout and class: the article's own class is \texttt{amsart} with packages including dsfont, amsxtra, amsmath, amscd, amssymb, amsfonts, mathrsfs, ytableau, amsthm, color, upgreek, tikz, float, tikz-cd; the wrapper is the lane's pinned \texttt{amsart} editorial wrapper at 10pt on A4, which restates the theorem-like environments the retained text needs and adds the article's own macro definitions that the retained text needs (\texttt{\textbackslash min}), with extra wrapper packages (bm, bbm, dsfont, xspace, cancel, mathtools). The delivered numbering is the unit's own. Assets: no retained span contains a figure environment, a graphics-inclusion command, a PDF-page inclusion, a table or a TikZ picture; no image file is copied into this package, the package declares no asset dependency and no image file is redistributed with it. Licence: version arXiv:2608.17709v3 of this article is distributed under the Creative Commons Attribution 4.0 International licence, as recorded in the arXiv licence metadata for this exact version and shown on the abstract page https://arxiv.org/abs/2608.17709v3. A full rights notice is delivered at licenses/arxiv-2608.17709-CC-BY-4.0.txt. Characters: every retained excerpt is pure ASCII, so the delivered engine, XeLaTeX with its default Latin Modern text font, reports no missing character.
\begin{theorem}\label{Prop:TypeBC}
The double cosets $K_{r',d'}\backslash W_n/K_{r,d}$ are parametrized by pairs $(a,t)$ satisfying
\[
\max(0,r+r'-n)\le a\le \min(r,r'),
\qquad
m:=r+r'-a\le t\le n.
\]
For such a pair, a convenient representative is
\[
w_{a,t}(i)=
\begin{cases}
 i, & 1\le i\le a,\\
 r'+i-a, & a+1\le i\le r,\\
 a+i-r, & r+1\le i\le m,\\
 -(t+m+1-i), & m+1\le i\le t,\\
 i, & t+1\le i\le n.
\end{cases}
\]
Moreover,
\[
K_{r,d}\cap w_{a,t}^{-1}K_{r',d'}w_{a,t}
=
W_{[1,a]}
\times\mathfrak S_{[a+1,r]}
\times\mathfrak S_{[r+1,m]}
\times\mathfrak S_{[m+1,t]}
\times\mathfrak S_{[t+1,n]}.
\]
\end{theorem}

\begin{proof}
Put
\[
A=[1,r],\qquad B=[r+1,n],\qquad
C=[1,r'],\qquad D=[r'+1,n].
\]
We first classify the double cosets.  For $w\in W_n$, the right coset
$wK_{r,d}$ is completely determined by the signed $d$-subset
\[
S_w:=w(B)\subset\{\pm1,\ldots,\pm n\},
\]
where no two elements have the same absolute value.  Indeed,
$W_{[1,r]}$ acts only on the complementary coordinates, while
$\mathfrak S_{[r+1,n]}$ merely permutes the elements of $B$; conversely,
if $S_w=S_{w'}$, then $w^{-1}w'\in K_{r,d}$.

Take orbits under the left action of
$K_{r',d'}=W_C\times\mathfrak S_D$.  The factor $W_C$ may arbitrarily
permute the coordinates with absolute value in $C$ and change their
signs, whereas $\mathfrak S_D$ may permute the coordinates in $D$ but
preserves their signs.  Hence the orbit of $S_w$ is determined by the two
integers
\[
b:=\#\{x\in S_w:|x|\le r'\},
\qquad
s:=\#\{x\in S_w:x<0,\ |x|>r'\}.
\]
These invariants are complete: every orbit contains the unique canonical
signed subset
\[
\{a+1,\ldots,r'\}
\sqcup
\{-(m+1),\ldots,-t\}
\sqcup
\{t+1,\ldots,n\},
\]
where
\[
a:=r'-b,\qquad
m:=r+r'-a,\qquad
t:=m+s.
\]
Since $S_w$ has $d$ elements while $C$ and $D$ have cardinalities
$r'$ and $d'$, respectively, one has
\[
\max(0,d-d')\le b\le\min(d,r').
\]
Substituting $a=r'-b$ gives
\[
\max(0,r+r'-n)\le a\le\min(r,r').
\]
Moreover,
\[
d-b=n-r-r'+a=n-m,
\]
so $0\le s\le d-b$ is equivalent to $m\le t\le n$.
The displayed element $w_{a,t}$ sends $B$ precisely to this canonical
signed subset.  Thus the pairs $(a,t)$ give a complete and nonredundant
set of double-coset parameters.

It remains to compute the intersection subgroup.  Decompose
\[
[1,n]=I_1\sqcup I_2\sqcup I_3\sqcup I_4\sqcup I_5
\]
with
\[
I_1=[1,a],\quad I_2=[a+1,r],\quad I_3=[r+1,m],
\quad I_4=[m+1,t],\quad I_5=[t+1,n].
\]
By construction, $w_{a,t}$ maps $I_1$ and $I_3$ into $C$, and maps
$I_2,I_4,I_5$ into $D$; the only negative signs occur on $I_4$.
Therefore an element of $K_{r,d}$ lying in
$w_{a,t}^{-1}K_{r',d'}w_{a,t}$ must preserve the five intervals above.
Indeed, the two parabolics already force preservation of the common
refinement of the partitions $A\sqcup B$ and
$w_{a,t}^{-1}(C)\sqcup w_{a,t}^{-1}(D)$, while a permutation mixing
$I_4$ with $I_5$ would, after conjugation by $w_{a,t}$, introduce a sign
change in the unsigned $D$-block.

Finally, signs are allowed exactly on $I_1$: the right parabolic permits
sign changes on $A=I_1\sqcup I_2$, whereas the conjugated left parabolic
permits them on $w_{a,t}^{-1}(C)=I_1\sqcup I_3$.  Their intersection
therefore permits arbitrary signed permutations on $I_1$ and only
ordinary permutations on the remaining four blocks.  Hence
\[
K_{r,d}\cap w_{a,t}^{-1}K_{r',d'}w_{a,t}
=
W_{I_1}\times\mathfrak S_{I_2}\times\mathfrak S_{I_3}
\times\mathfrak S_{I_4}\times\mathfrak S_{I_5},
\]
which is the asserted formula.
\end{proof}

\end{document}
